% GENERATED FILE. Edit canonical lessons, then run npm run generate:books.
% canonical-source-hash: fnv1a64:075ab71f1d505df6
% canonical-lessons: KA-C504-siduku, KA-C504-matugara, KA-C504-mungopi, KA-C504-hedi, KA-C504-todalu, KA-R504-review-498-504

\chapter{Talk and Temper}
\label{ch:ka-talk-and-temper}
\hlchaptermodality{504}

\begin{chapteropening}
\textbf{By the end of this chapter:} \emph{I can say to snap irritably, a talker, short-tempered, a coward and a stammer.}

Complete the last lesson of chapter 504: I can say to snap irritably, a talker, short-tempered, a coward and a stammer.
\end{chapteropening}

\section[\texorpdfstring{\kn{ಸಿಡುಕು} (siḍuku)}{ಸಿಡುಕು (siḍuku)}]{\kn{ಸಿಡುಕು} (siḍuku) — to snap irritably}

\label{lesson:KA-C504-siduku}



\begin{quote}
Before the new one: say the Kannada for forgetfulness, then the Kannada for concentration.
\end{quote}

\subsection*{You'll want to know: \kn{ಸಿಡುಕು}}
\textbf{\kn{ಸಿಡುಕು}} — \emph{siḍuku} — \textquotedblleft{}to snap irritably\textquotedblright{}.

\begin{quote}
\kn{ಹಸಿವಾದಾಗ} \kn{ನನ್ನ} \kn{ತಮ್ಮ} \kn{ಎಲ್ಲರ} \kn{ಮೇಲೆ} \kn{ಸಿಡುಕುತ್ತಾನೆ}.
\end{quote}

\begin{quote}
\emph{hasivādāga nanna tamma ellara mēle siḍukuttāne.}
\end{quote}

\begin{quote}
When he is hungry, my younger brother snaps at everyone.
\end{quote}

Past (I …): \kn{ಸಿಡುಕಿದೆ}, \emph{siḍukide}.

At whom: … \kn{ಮೇಲೆ} \kn{ಸಿಡುಕು}. Also a noun, irritability: \kn{ಸಿಡುಕಿನ} \kn{ಮುಖ}, a sour face.

\begin{cousinweb}
\textbf{\kn{ಸಿಡುಕು}} (\emph{siḍuku}) is a native Kannada verb, \textquotedblleft{}to be irritable, to snap\textquotedblright{}; its origin is not settled.
\end{cousinweb}

\begin{grammarlens}[title={What you've built}]
One more word for describing what you lived and felt.
\end{grammarlens}

\subsection*{Your turn}
\begin{itemize}
\raggedright
  \item \emph{Say it:} \emph{siḍuku}
  \item \emph{Say it:} \emph{siḍuku}, once more
  \item \emph{Say it:} \emph{hasivādāga nanna tamma ellara mēle siḍukuttāne.} — When he is hungry, my younger brother snaps at everyone.
  \item \emph{Recall:} say the Kannada for forgetfulness, then the Kannada for concentration, then say \emph{siḍuku} again
\end{itemize}

\begin{tcolorbox}[breakable,colback=teal!4,colframe=teal!35!black,title={Before you move on}]
What does \kn{ಸಿಡುಕು} mean? (\textquotedblleft{}To snap irritably\textquotedblright{}.) Say it once more.
\end{tcolorbox}
\section[\texorpdfstring{\kn{ಮಾತುಗಾರ} (mātugāra)}{ಮಾತುಗಾರ (mātugāra)}]{\kn{ಮಾತುಗಾರ} (mātugāra) — a talker; a good speaker}

\label{lesson:KA-C504-matugara}



\begin{quote}
Before the new one: say the Kannada for concentration, then the Kannada for to snap irritably.
\end{quote}

\subsection*{You'll want to know: \kn{ಮಾತುಗಾರ}}
\textbf{\kn{ಮಾತುಗಾರ}} — \emph{mātugāra} — \textquotedblleft{}a talker; a good speaker\textquotedblright{}.

\begin{quote}
\kn{ನನ್ನ} \kn{ಮಾವ} \kn{ದೊಡ್ಡ} \kn{ಮಾತುಗಾರ}.
\end{quote}

\begin{quote}
\emph{nanna māva doḍḍa mātugāra.}
\end{quote}

\begin{quote}
My uncle is a great talker.
\end{quote}

\kn{ಮಾತು}, speech + -\kn{ಗಾರ}, \textquotedblleft{}one who does it\textquotedblright{}, as in \kn{ಆಟಗಾರ}.

\begin{cousinweb}
\textbf{\kn{ಮಾತುಗಾರ}} (\emph{mātugāra}) is native Dravidian: DEDR 4834 lists Kannada \emph{mātu}, \textquotedblleft{}a word, speech\textquotedblright{}, and \emph{mātugāṟa}, \textquotedblleft{}a talkative man\textquotedblright{}, with the suffix -\emph{gāra}, \textquotedblleft{}one who does\textquotedblright{}, as in \kn{ಆಟಗಾರ}.
\end{cousinweb}

\begin{grammarlens}[title={What you've built}]
One more word for describing what you lived and felt.
\end{grammarlens}

\subsection*{Your turn}
\begin{itemize}
\raggedright
  \item \emph{Say it:} \emph{mātugāra}
  \item \emph{Say it:} \emph{mātugāra}, once more
  \item \emph{Say it:} \emph{nanna māva doḍḍa mātugāra.} — My uncle is a great talker.
  \item \emph{Recall:} say the Kannada for concentration, then the Kannada for to snap irritably, then say \emph{mātugāra} again
\end{itemize}

\begin{tcolorbox}[breakable,colback=teal!4,colframe=teal!35!black,title={Before you move on}]
What does \kn{ಮಾತುಗಾರ} mean? (\textquotedblleft{}A talker; a good speaker\textquotedblright{}.) Say it once more.
\end{tcolorbox}
\section[\texorpdfstring{\kn{ಮುಂಗೋಪಿ} (muṅgōpi)}{ಮುಂಗೋಪಿ (muṅgōpi)}]{\kn{ಮುಂಗೋಪಿ} (muṅgōpi) — short-tempered}

\label{lesson:KA-C504-mungopi}



\begin{quote}
Before the new one: say the Kannada for to snap irritably, then the Kannada for a talker.
\end{quote}

\subsection*{You'll want to know: \kn{ಮುಂಗೋಪಿ}}
\textbf{\kn{ಮುಂಗೋಪಿ}} — \emph{muṅgōpi} — \textquotedblleft{}short-tempered\textquotedblright{}.

\begin{quote}
\kn{ನಮ್ಮ} \kn{ಮೇಲಧಿಕಾರಿ} \kn{ಮುಂಗೋಪಿ}, \kn{ಬೇಗ} \kn{ಕೋಪಗೊಳ್ಳುತ್ತಾರೆ}.
\end{quote}

\begin{quote}
\emph{namma mēladhikāri muṅgōpi, bēga kōpagoḷḷuttāre.}
\end{quote}

\begin{quote}
Our boss is short-tempered; he gets angry quickly.
\end{quote}

\kn{ಮುನ್}, before + \kn{ಕೋಪ}, anger: angry before anything has happened.

\begin{cousinweb}
\textbf{\kn{ಮುಂಗೋಪಿ}} (\emph{muṅgōpi}) is a Kannada compound: native \emph{mun}, \textquotedblleft{}before, in front\textquotedblright{} (DEDR 5020), and \emph{kōpa}, \textquotedblleft{}anger\textquotedblright{}, borrowed from Sanskrit (CDIAL 3516), with -\emph{i} for the person.
\end{cousinweb}

\begin{grammarlens}[title={What you've built}]
One more word for describing what you lived and felt.
\end{grammarlens}

\subsection*{Your turn}
\begin{itemize}
\raggedright
  \item \emph{Say it:} \emph{muṅgōpi}
  \item \emph{Say it:} \emph{muṅgōpi}, once more
  \item \emph{Say it:} \emph{namma mēladhikāri muṅgōpi, bēga kōpagoḷḷuttāre.} — Our boss is short-tempered; he gets angry quickly.
  \item \emph{Recall:} say the Kannada for to snap irritably, then the Kannada for a talker, then say \emph{muṅgōpi} again
\end{itemize}

\begin{tcolorbox}[breakable,colback=teal!4,colframe=teal!35!black,title={Before you move on}]
What does \kn{ಮುಂಗೋಪಿ} mean? (\textquotedblleft{}Short-tempered\textquotedblright{}.) Say it once more.
\end{tcolorbox}
\section[\texorpdfstring{\kn{ಹೇಡಿ} (hēḍi)}{ಹೇಡಿ (hēḍi)}]{\kn{ಹೇಡಿ} (hēḍi) — a coward}

\label{lesson:KA-C504-hedi}



\begin{quote}
Before the new one: say the Kannada for a talker, then the Kannada for short-tempered.
\end{quote}

\subsection*{You'll want to know: \kn{ಹೇಡಿ}}
\textbf{\kn{ಹೇಡಿ}} — \emph{hēḍi} — \textquotedblleft{}a coward\textquotedblright{}.

\begin{quote}
\kn{ಹೇಡಿ} \kn{ಆಗಬೇಡ} \kn{ಅಂತ} \kn{ಅಪ್ಪ} \kn{ಹೇಳಿದರು}.
\end{quote}

\begin{quote}
\emph{hēḍi āgabēḍa anta appa hēḷidaru.}
\end{quote}

\begin{quote}
Father said, do not be a coward.
\end{quote}

The opposite is \kn{ಧೈರ್ಯವಂತ}, brave.

\begin{cousinweb}
\textbf{\kn{ಹೇಡಿ}} (\emph{hēḍi}) is native Dravidian: DEDR 4434 sets Kannada \emph{hēḍi} beside Tamil \emph{pēṭi}, where fear is one of the senses.
\end{cousinweb}

\begin{grammarlens}[title={What you've built}]
One more word for describing what you lived and felt.
\end{grammarlens}

\subsection*{Your turn}
\begin{itemize}
\raggedright
  \item \emph{Say it:} \emph{hēḍi}
  \item \emph{Say it:} \emph{hēḍi}, once more
  \item \emph{Say it:} \emph{hēḍi āgabēḍa anta appa hēḷidaru.} — Father said, do not be a coward.
  \item \emph{Recall:} say the Kannada for a talker, then the Kannada for short-tempered, then say \emph{hēḍi} again
\end{itemize}

\begin{tcolorbox}[breakable,colback=teal!4,colframe=teal!35!black,title={Before you move on}]
What does \kn{ಹೇಡಿ} mean? (\textquotedblleft{}A coward\textquotedblright{}.) Say it once more.
\end{tcolorbox}
\section[\texorpdfstring{\kn{ತೊದಲು} (todalu)}{ತೊದಲು (todalu)}]{\kn{ತೊದಲು} (todalu) — a stammer; baby talk}

\label{lesson:KA-C504-todalu}



\begin{quote}
Before the new one: say the Kannada for short-tempered, then the Kannada for a coward.
\end{quote}

\subsection*{You'll want to know: \kn{ತೊದಲು}}
\textbf{\kn{ತೊದಲು}} — \emph{todalu} — \textquotedblleft{}a stammer; baby talk\textquotedblright{}.

\begin{quote}
\kn{ಮಗುವಿನ} \kn{ತೊದಲು} \kn{ಮಾತು} \kn{ಎಲ್ಲರಿಗೂ} \kn{ಇಷ್ಟ}.
\end{quote}

\begin{quote}
\emph{maguvina todalu mātu ellarigū iṣṭa.}
\end{quote}

\begin{quote}
Everyone loves the baby's lisping talk.
\end{quote}

\kn{ತೊದಲು} \kn{ಮಾತು}, baby talk. It is a verb as well: \kn{ತೊದಲುತ್ತಾ} \kn{ಮಾತಾಡು}, to speak stammering.

\begin{cousinweb}
\textbf{\kn{ತೊದಲು}} (\emph{todalu}) is native Dravidian: DEDR 3503 lists Kannada \emph{todal(u)}, \textquotedblleft{}faltering, stammering, speaking indistinctly\textquotedblright{}.
\end{cousinweb}

\begin{grammarlens}[title={What you've built}]
One more word for describing what you lived and felt.
\end{grammarlens}

\subsection*{Your turn}
\begin{itemize}
\raggedright
  \item \emph{Say it:} \emph{todalu}
  \item \emph{Say it:} \emph{todalu}, once more
  \item \emph{Say it:} \emph{maguvina todalu mātu ellarigū iṣṭa.} — Everyone loves the baby's lisping talk.
  \item \emph{Recall:} say the Kannada for short-tempered, then the Kannada for a coward, then say \emph{todalu} again
\end{itemize}

\begin{tcolorbox}[breakable,colback=teal!4,colframe=teal!35!black,title={Before you move on}]
What does \kn{ತೊದಲು} mean? (\textquotedblleft{}A stammer; baby talk\textquotedblright{}.) Say it once more.
\end{tcolorbox}
\section[(dialogue)]{Looking back: from Trust and Sympathy to Talk and Temper}

\label{lesson:KA-R504-review-498-504}



\begin{quote}
Say the Kannada for to believe and to forgive.
\end{quote}

\subsection*{Your turn}
Say these aloud:

\begin{itemize}
\raggedright
  \item to believe, to forgive, suspicion, compassion, sympathy
  \item to sulk, to smile, to hesitate, envy, self-confidence
  \item temperament, personality, conduct, attitude, to imitate
  \item gentle, innocent, cruel, harshly frank, to mimic mockingly
  \item polite, obedient, discipline, to take note, to react
  \item quick, talent, slow-witted, forgetfulness, concentration
  \item to snap irritably, a talker, short-tempered, a coward, a stammer
\end{itemize}

\begin{tcolorbox}[breakable,colback=teal!4,colframe=teal!35!black,title={Before you move on}]
To believe? (\textbf{\kn{ನಂಬು}}, \emph{nambu}.) To forgive? (\textbf{\kn{ಮನ್ನಿಸು}}, \emph{mannisu}.) Suspicion? (\textbf{\kn{ಸಂಶಯ}}, \emph{saṃśaya}.) Compassion? (\textbf{\kn{ಕರುಣೆ}}, \emph{karuṇe}.) Sympathy? (\textbf{\kn{ಸಹಾನುಭೂತಿ}}, \emph{sahānubhūti}.) To sulk? (\textbf{\kn{ಮುನಿಸಿಕೊ}}, \emph{munisiko}.) To smile? (\textbf{\kn{ಮುಗುಳ್ನಗು}}, \emph{muguḷnagu}.) To hesitate? (\textbf{\kn{ಹಿಂಜರಿ}}, \emph{hiñjari}.) Envy? (\textbf{\kn{ಹೊಟ್ಟೆಕಿಚ್ಚು}}, \emph{hoṭṭekiccu}.) Self-confidence? (\textbf{\kn{ಆತ್ಮವಿಶ್ವಾಸ}}, \emph{ātmaviśvāsa}.) Temperament? (\textbf{\kn{ಸ್ವಭಾವ}}, \emph{svabhāva}.) Personality? (\textbf{\kn{ವ್ಯಕ್ತಿತ್ವ}}, \emph{vyaktitva}.) Conduct? (\textbf{\kn{ನಡತೆ}}, \emph{naḍate}.) Attitude? (\textbf{\kn{ವರ್ತನೆ}}, \emph{vartane}.) To imitate? (\textbf{\kn{ಅನುಕರಿಸು}}, \emph{anukarisu}.) Gentle? (\textbf{\kn{ಸೌಮ್ಯ}}, \emph{saumya}.) Innocent? (\textbf{\kn{ಮುಗ್ಧ}}, \emph{mugdha}.) Cruel? (\textbf{\kn{ಕ್ರೂರ}}, \emph{krūra}.) Harshly frank? (\textbf{\kn{ನಿಷ್ಠುರ}}, \emph{niṣṭhura}.) To mimic mockingly? (\textbf{\kn{ಅಣಕಿಸು}}, \emph{aṇakisu}.) Polite? (\textbf{\kn{ಸಭ್ಯ}}, \emph{sabhya}.) Obedient? (\textbf{\kn{ವಿಧೇಯ}}, \emph{vidhēya}.) Discipline? (\textbf{\kn{ಶಿಸ್ತು}}, \emph{śistu}.) To take note? (\textbf{\kn{ಗಮನಿಸು}}, \emph{gamanisu}.) To react? (\textbf{\kn{ಸ್ಪಂದಿಸು}}, \emph{spandisu}.) Quick? (\textbf{\kn{ಚುರುಕು}}, \emph{curuku}.) Talent? (\textbf{\kn{ಪ್ರತಿಭೆ}}, \emph{pratibhe}.) Slow-witted? (\textbf{\kn{ದಡ್ಡ}}, \emph{daḍḍa}.) Forgetfulness? (\textbf{\kn{ಮರೆವು}}, \emph{marevu}.) Concentration? (\textbf{\kn{ಏಕಾಗ್ರತೆ}}, \emph{ēkāgrate}.) To snap irritably? (\textbf{\kn{ಸಿಡುಕು}}, \emph{siḍuku}.) A talker? (\textbf{\kn{ಮಾತುಗಾರ}}, \emph{mātugāra}.) Short-tempered? (\textbf{\kn{ಮುಂಗೋಪಿ}}, \emph{muṅgōpi}.) A coward? (\textbf{\kn{ಹೇಡಿ}}, \emph{hēḍi}.) A stammer? (\textbf{\kn{ತೊದಲು}}, \emph{todalu}.)
\end{tcolorbox}
